\documentclass[12pt]{article}
\usepackage[margin=2in]{geometry}
\usepackage{amsmath,amsthm,amssymb,mathtools}
\usepackage{float}
\usepackage{setspace}
\setstretch{1.0}
\usepackage{hyperref}

\newtheorem*{theorem}{Theorem}
\newcommand{\N}{\mathbb{N}}

\begin{document}

\section{Introduction to Lean}

\begin{itemize}
  \item Functional programming language and theorem prover
  \item Designed for formal verification tasks
  \item Based on dependent type theory
  \item Open-source, first developed by Microsoft
  \item First launched in 2013; Lean 4 released in 2021
  \item In 2026, three Erd{\H o}s conjectures were resolved in Lean
\end{itemize}

\section{Mathlib}
\begin{itemize}
  \item Library (or package) for Lean
  \item Repository of formalised mathematics
  \item Started in 2017; now contains over 200\,000 theorems
\end{itemize}

\pagebreak

\section{Writing our first proofs}

\pagebreak

\section{Review of syntax}
\begin{itemize}
  \item \texttt{example} \\
    Used to introduce a new proposition to prove.
  \item \texttt{rfl} \\
    Prove a proposition by definition.
  \item \texttt{exact h} \\
    Prove a proposition by applying hypothesis \texttt{h}.
  \item \texttt{rw [h]} \\
    Rewrite the goal using hypothesis \texttt{h}.
  \item Comments can be written with: \texttt{-- here is a comment}.
  \item \texttt{\#check} \\
    Check the type of an expression. \\
    \texttt{\#check 1} gives \texttt{1:}$\mathbb{N}$. \\
    \texttt{\#check (1 + 1 = 2)} gives \texttt{(1 + 1 = 2):Prop}.
  \item \texttt{constructor} \\
    Transform the goal \texttt{P and Q} into two goals, \texttt{P} and \texttt{Q}.
  \item \texttt{sorry} \\
    Temporarily prove any goal (and raise a warning).
  \item \texttt{have} \\
    Introduce a new sub-goal to prove.
\end{itemize}

\pagebreak

\section{Basics of type-theoretic logic}

\begin{itemize}
  \item The Curry--Howard correspondence gives an isomorphism between
    models for \emph{type theory} and \emph{logic}.
\end{itemize}

\begin{table}[H]
  \begin{center}
    \begin{tabular}{|c|c|}
      \hline
      Type theory & Logic \\
      \hline
      Type \texttt{T} & Proposition \texttt{P} \\
      Element of type \texttt{x:T} & Proof of proposition \texttt{h:P} \\
      \texttt{n:}$\mathbb{N}$ & \texttt{rfl:(1 + 1 = 2)} \\
      \hline
    \end{tabular}
  \end{center}
\end{table}

\begin{itemize}
  \item The proposition \texttt{True} has a proof
    \texttt{trivial:True}.
  \item The proposition \texttt{False} has no proofs!
  \item If \texttt{P} and, \texttt{Q} are propositions,
    then a proof of \texttt{P}$\implies$\texttt{Q}
    is a function mapping proofs of \texttt{P}
    to proofs of \texttt{Q}.
  \item A proof of \texttt{P and Q} is a pair of proofs,
    one for \texttt{P} and one for \texttt{Q}:
    \texttt{(hP, hQ):(P and Q)} where
    \texttt{hP:P} and \texttt{hQ:Q}.
  \item A proof of \texttt{not P} is a function
    mapping proofs of \texttt{P} to proofs of \texttt{False}.
  \item \texttt{P or Q} is rather more complicated:
    we may prove \texttt{P or Q} without having a proof for either
    \texttt{P} or \texttt{Q}. For example:
    %
    \begin{align*}
      \exists p, q \in \mathbb{R} \setminus \mathbb{Q}:
      p^q \in \mathbb{Q}.
    \end{align*}
    \emph{Proof.}
    Consider $x \coloneqq \sqrt{2}^{\sqrt{2}}$ and
    $x^{\sqrt 2}$.
\end{itemize}

\pagebreak

\section{Infinitude of the primes}

We present a famous result and proof due to Euclid,
published in the Elements, around 300 BCE.

\begin{theorem}[Infinitude of the primes]
  For all $n \in \N$,
  there exists $p \in \N$ such that
  $p > n$ and $p$ is prime.
\end{theorem}

\begin{proof}
  %
  Take $n \in \N$, write $m \coloneqq n! + 1$,
  and let $p$ be a prime factor of $m$.
  %
  \begin{itemize}
    %
    \item Note that $m \neq 1$ because $n! \neq 0$.
    \item Therefore, $p$ is prime.
    \item Suppose for contradiction that we do not have $p > n$.
    \item Then $p \leq n$.
    \item Note that $p > 0$ because $p$ is a factor of $m$.
    \item We have $p \mid m$ by definition of $p$.
    \item Also, $p \mid n!$ because $p \leq n$ by assumption and $p > 0$.
    \item So $p \mid 1$ by modular arithmetic.
    \item Thus, $p = 1$.
    \item But $1$ is not prime \textreferencemark.
      \qedhere
    %
  \end{itemize}
  %
\end{proof}

\pagebreak

\section{Getting help}

\begin{itemize}
  \item \texttt{\#check} \\
    Check the type of an expression.
    Remember: propositions and proofs have types
    and can be checked! \\
    For example, \texttt{\#check Nat.eq\_one\_of\_dvd\_one}
  \item \texttt{\#eval} \\
    Evaluate some Lean code. \\
    For example, \texttt{\#eval Nat.minFac 1}
  \item Advanced tactics for finding proofs: \\
    \texttt{aesop}, \texttt{omega}, \texttt{linarith},
    \texttt{ring}, \texttt{grind}
  \item \url{https://leansearch.net/} \\
    Search Mathlib using natural language.
  \item \url{https://github.com/leanprover-community/mathlib4} \\
    The Mathlib code repository.
  \item https://lean-lang.org/ \\
    Core documentation for Lean.
\end{itemize}

\pagebreak

\section{A simple limit theorem}

\end{document}
